% CN: 共享侧响应稳定性及其与势传播的联系；辅助预测分析已归档。
% EN: Shared-side response stability and its link to potential propagation; auxiliary prediction analysis is archived.
\subsection{Inter-Level Response Stability}
\label{app:uot-hierarchy-stability}

Consider two KL-UOT problems sharing target locations, positive reference weights $\mathbf b$, and penalty $\rho_b>0$, while their source locations, positive source weights, or positive source penalties may differ.
Costs are finite, and the coupling has no entropic regularization.
For a target potential $\mathbf g$, define its normalized marginal response by
\begin{equation*}
p_j(\mathbf g)=\frac{b_j e^{-g_j/\rho_b}}{\sum_k b_k e^{-g_k/\rho_b}}.
\end{equation*}
This response is invariant under scalar translation of $\mathbf g$.
We measure response differences by $H^2(p,z)=\sum_j(\sqrt{p_j}-\sqrt{z_j})^2$.
The argument applies to adjacent hierarchy levels sharing the target side; exchanging the marginals covers the other direction.

For two source instances with a common target set, take their union as the source index set and define the rectangular cost oscillation
$$\Gamma=\max_{i,i',j,j'}|C_{ij}-C_{ij'}-C_{i'j}+C_{i'j'}|.$$

\begin{proposition}[Shared-side response stability]
\label{prop:uot-hierarchy-stability}
\label{prop:uot-hierarchy-interaction}
Let $p^0,p^1$ be the normalized optimal target marginals of any two such source instances with the same $b$ and $\rho_b$.
Then
$$H^2(p^0,p^1)\leq2\left[1-\operatorname{sech}\!\left(\frac{\Gamma}{4\rho_b}\right)\right]\leq\frac{\Gamma^2}{16\rho_b^2}.$$
For squared Euclidean cost, $\Gamma\leq2\operatorname{diam}(S)\operatorname{diam}(T)$, where $S$ contains both source sets and $T$ is the common target set.
For separable cost $C_{ij}=u_i+v_j$, the target-response shape is exactly $p_j\propto b_j e^{-v_j/\rho_b}$ and is unchanged by source subsampling or source reweighting.
The first bound is sharp over pairs of source instances.
\end{proposition}
\begin{proof}
Positive optimal target marginals imply $g_j=\min_i(C_{ij}-f_i)$.
Thus each difference $g_j-g_{j'}$ lies between the minimum and maximum of $C_{ij}-C_{ij'}$ over the available source indices.
Using the union of both source sets gives $\operatorname{osc}(g^1-g^0)\leq\Gamma$.
The normalized log-likelihood ratio therefore has oscillation at most $D=\Gamma/\rho_b$.

For completeness, let $x_j=\sqrt{p^1_j/p^0_j}\in[\alpha,\beta]$, with $\beta/\alpha\leq e^{D/2}$.
Since $\mathbb E_{p^0}x^2=1$ and $(x-\alpha)(x-\beta)\leq0$,
$$\mathbb E_{p^0}x\geq\frac{1+\alpha\beta}{\alpha+\beta}\geq\operatorname{sech}(D/4).$$
The final inequality follows by minimizing over $\alpha$ at fixed $\beta/\alpha$.
The affinity identity $H^2=2-2\mathbb E x$ proves the first bound.
The inequality $2(1-\operatorname{sech} t)\leq t^2$ proves the second.
For quadratic cost the rectangular difference is $-2(s_i-s_{i'})^\top(t_j-t_{j'})$.
If $\Gamma=0$, the bound gives exact inheritance; the displayed separable response also follows directly from the optimality conditions.
For sharpness, take two singleton source instances whose costs to two equally weighted targets are $(0,\Gamma/2)$ and $(\Gamma/2,0)$.
Their normalized target responses attain $H^2=2[1-\operatorname{sech}(\Gamma/(4\rho_b))]$.
\end{proof}

The bound concerns optimal response shapes and can be conservative on large supports; it provides neither a sampling rate nor a warm-start iteration bound.

\begin{corollary}[Weighted potential stability]
\label{cor:uot-hierarchy-weighted-potentials}
For any two finite target potentials $g^0,g^1$ with a common positive reference $b$ and penalty $\rho_b>0$, define
$$Z_s=\sum_j b_j e^{-g_j^s/\rho_b},\qquad p_j^s=\frac{b_j e^{-g_j^s/\rho_b}}{Z_s},\qquad s\in\{0,1\},$$
and remove scalar translations by setting $\bar g^s=g^s+\rho_b\log Z_s\,\mathbf1$.
Then
\begin{equation*}
\sum_j\sqrt{p_j^0p_j^1}\,(\bar g_j^1-\bar g_j^0)^2\leq4\rho_b^2H^2(p^0,p^1).
\end{equation*}
No uniform positive lower bound on the response entries is required.
\end{corollary}
\begin{proof}
Writing $\delta_j=\bar g_j^1-\bar g_j^0$, the response identity gives $\delta_j=-\rho_b\log(p_j^1/p_j^0)$.
Consequently,
$$H^2(p^0,p^1)=4\sum_j\sqrt{p_j^0p_j^1}\,\sinh^2\!\left(\frac{\delta_j}{4\rho_b}\right)\geq\frac{1}{4\rho_b^2}\sum_j\sqrt{p_j^0p_j^1}\,\delta_j^2,$$
where the inequality uses $|\sinh x|\geq|x|$.
\end{proof}

The estimate connects response stability to dual propagation by controlling the aligned potentials with weights that emphasize coordinates carrying mass in both responses.
Coordinates with vanishing response weight may still have large potential differences, so the estimate does not imply uniform coordinatewise stability.
